\section{Conclusion}
\label{sec:conclusion}

In this work, we explored the application of generative imitation learning for bimanual garment folding in the context of the ICRA 2026 LeHome Challenge. Our investigation led us to shift from our initial Diffusion Policy proposal to an architecture based on Action Chunking with Transformers (ACT). We demonstrated that ACT's ability to model action sequences as chunks is significantly better suited for coordinated bimanual trajectories compared to standard Diffusion Policy or large-scale VLA models when trained on specialized datasets.

By integrating geometric priors through colorized depth and employing a ResNet-18 visual backbone, we achieved a macro-average success rate of 62.92\% and demonstrated superior generalization (51.25\%) on unseen garments. While higher-capacity backbones like DINOv2 achieved slightly higher peaks on seen data (63.54\%), the ResNet-18 configuration provided a more robust balance between performance and generalization, serving as our final submission to the challenge. Our findings highlight the importance of choosing architectures that capture complex temporal dependencies and the effectiveness of mid-level geometric representations in the challenging domain of deformable object manipulation.
